% Приватне дослідження: лише агреговані операційні перевірки.
\section{Польове дослідження на архіві автора}\label{sec:field}

Приватне дослідження запускає реалізований адаптивний пошук на наявному архіві автора з трьома наборами запитів, написаних під час розроблення системи. Їхні назви (main, holdout7 і holdout11) --- позначення розроблення, а не свідчення незалежної вибірки чи проспективного користувацького дослідження. Набори містять відповідно 80, 86 і 214 запитів. Документи, тексти запитів, шляхи, фрагменти й індивідуальні результати не експортуються.

Виконуються три операційні перевірки. Для точного запиту успіх означає наявність очікуваної фрази у знайденому фрагменті після нормалізації пробілів і регістру; ранги й обернені ранги стосуються фрагментів. Для описового запиту успіх означає наявність очікуваного проєкту серед перших п’яти фрагментів. Для запиту без відповіді успіх означає відсутність результатів у лексичному режимі; це не перевірка утримання від відповіді векторного пошуку. Ці перевірки слабші за незалежні оцінки релевантності документів. Кожен набір і тип підсумовуються окремо з 95\% бутстреп-інтервалами за запитами; набори не об’єднуються в загальну оцінку ефективності.

\ifpaperfield
\begin{table}[!htbp]
\centering
\caption{Операційні перевірки приватного архіву: кількості, середні та 95\% бутстреп-інтервали за запитами. Для точних запитів ранги стосуються фрагментів; описові перевіряють проєкт; безвідповідні --- лексичне утримання. Експортуються лише агрегати.}
\label{tab:field}
% Generated from primary=gemini_reviewed; see ../manifest.json for source SHA-256.
Архів: 11\,394 документів, 902\,739 фрагментів.
\par\smallskip
% private archive aggregates only; exact ranks are chunk ranks; source field_aggregate.json
\begingroup\small\setlength{\tabcolsep}{4pt}\renewcommand{\arraystretch}{1.3}
\noindent\begin{tabularx}{\linewidth}{@{}l*{4}{>{\centering\arraybackslash}X}@{}}
\toprule
Набір & Тип & Запитів & Метрика & Середнє [95\% ДІ] \\
\midrule
holdout11 & exact & 192 & hit@1 & 0.672 [0.604, 0.740] \\
holdout11 & exact & 192 & hit@5 & 0.708 [0.641, 0.771] \\
holdout11 & exact & 192 & mrr@10 & 0.691 [0.627, 0.754] \\
holdout11 & absent & 10 & pass & 1.000 [1.000, 1.000] \\
holdout11 & described & 12 & hit@5 & 1.000 [1.000, 1.000] \\
holdout7 & exact & 64 & hit@1 & 0.641 [0.516, 0.750] \\
holdout7 & exact & 64 & hit@5 & 0.703 [0.594, 0.812] \\
holdout7 & exact & 64 & mrr@10 & 0.668 [0.557, 0.777] \\
holdout7 & absent & 10 & pass & 1.000 [1.000, 1.000] \\
holdout7 & described & 12 & hit@5 & 1.000 [1.000, 1.000] \\
main & exact & 66 & hit@1 & 0.697 [0.576, 0.803] \\
main & exact & 66 & hit@5 & 0.727 [0.621, 0.833] \\
main & exact & 66 & mrr@10 & 0.709 [0.598, 0.813] \\
main & absent & 8 & pass & 1.000 [1.000, 1.000] \\
main & described & 6 & hit@5 & 1.000 [1.000, 1.000] \\
\bottomrule
\end{tabularx}\endgroup

\end{table}
Табл.~\ref{tab:field} описує це середовище й ці запити розроблення. Якщо всі спостереження успішні, перцентильний бутстреп дає вироджений інтервал [1,1]; це не гарантія для популяції. Інтервали характеризують перевибірку всередині набору; вони не враховують налаштування на цьому архіві, добір запитів чи відмінності між людьми.
\else
\pending{Агрегати польового дослідження ще не отримано}
\fi
